\section{Architecture et fonctionnement interne}

\subsection{Vue d’ensemble}

KORA suit une architecture locale à responsabilités séparées. La popup porte
l’intention de l’utilisateur; le content script observe la page; le moteur
décide si un segment peut être traduit; les modèles linguistiques classent
uniquement des formulations déjà autorisées. Le site cible demeure maître de
ses données et de sa logique métier.

\begin{figure}[H]
\centering
\resizebox{\linewidth}{!}{%
\begin{tikzpicture}[
  box/.style={draw=Line,fill=white,rounded corners=2mm,minimum height=12mm,
    minimum width=29mm,align=center,font=\small},
  core/.style={box,draw=Teal,fill=Mint},
  ext/.style={box,draw=InkSoft,fill=Paper},
  data/.style={draw=Gold,fill=Warm,rounded corners=2mm,minimum height=10mm,
    minimum width=28mm,align=center,font=\small},
  arrow/.style={-{Latex[length=2mm]},thick,color=TealDark},
  bi/.style={{Latex[length=2mm]}-{Latex[length=2mm]},thick,color=TealDark}
]
\node[ext] (user) {Utilisateur};
\node[core,right=13mm of user] (popup) {Popup\\activation · langue};
\node[data,below=7mm of popup] (storage) {Stockage local\\préférences seules};
\node[core,right=13mm of popup] (content) {Content script\\DOM · mutations};
\node[core,right=13mm of content] (engine) {Moteur\\sens · intégrité};
\node[data,above right=6mm and 10mm of engine] (mem) {Mémoires et\\gabarits FR/EN};
\node[data,right=10mm of engine] (spell) {Correcteur\\orthographique};
\node[data,below right=6mm and 10mm of engine] (model) {Modèles \code{fub}\\neuronal + n-gramme};
\node[ext,below=20mm of content] (page) {Page financière existante\\texte visible · formulaire · application};

\draw[bi] (user) -- (popup);
\draw[bi] (popup) -- node[above,font=\scriptsize]{messages Chrome} (content);
\draw[bi] (popup) -- (storage);
\draw[bi] (content) -- (storage);
\draw[bi] (content) -- node[above,font=\scriptsize]{segments} (engine);
\draw[bi] (content) -- node[left,font=\scriptsize]{lecture / rendu} (page);
\draw[bi] (engine) -- (mem);
\draw[bi] (engine) -- (spell);
\draw[bi] (engine) -- (model);
\end{tikzpicture}}
\caption{Architecture d’exécution : toutes les ressources linguistiques sont
embarquées dans l’extension.}
\end{figure}

\begin{evidencebox}[Absence de dépendance serveur]
Le bundle n’utilise aucune API de traduction, aucun proxy et aucun code
distant. La connexion Internet sert uniquement à charger le site financier
visité. Déconnecter le réseau après chargement ne change pas le pipeline KORA.
\end{evidencebox}

\subsection{Rôle précis de chaque composant}

\begin{longtable}{@{}|L{43mm}|L{58mm}|L{55mm}|@{}}
\hline
\textbf{Composant} & \textbf{Responsabilité} & \textbf{Contrat de sûreté}\\
\hline
\endfirsthead
\hline
\textbf{Composant} & \textbf{Responsabilité} & \textbf{Contrat de sûreté}\\
\hline
\endhead
\code{manifest.json} &
Déclare MV3, popup, icônes, content script et permissions. &
Pas de permission réseau ni de code distant.\\
\hline
\code{src/popup.*} &
Choix Fulfulde/Français/English, activation et état. &
Ne reçoit pas le contenu de page; rend le focus après action.\\
\hline
\code{src/content.js} &
Parcourt le DOM, collecte les valeurs, applique les sorties, restaure les
originaux et observe les mutations. &
Ne remplace pas les éléments interactifs; ne modifie aucune propriété de
formulaire.\\
\hline
\code{src/engine.js} &
Détection de source, cascade de traduction, couverture sémantique et contrôle
d’intégrité. &
Retourne l’original si le sens ou les valeurs échouent.\\
\hline
\code{src/spelling.js} &
Rapproche une entrée fautive d’une clé FR/EN connue avec distance
Damerau--Levenshtein bornée. &
Accepte seulement un meilleur sens unique; exclut nombres et codes.\\
\hline
\code{src/translations.js} &
Mémoire du parcours Safinum relue. &
Prioritaire sur les candidats terminologiques.\\
\hline
\texttt{src/general-\allowbreak translations.js} &
Libellés web et financiers généraux. &
Activé selon la langue et le contexte détectés.\\
\hline
\texttt{src/corpus-\allowbreak translations.js} &
Index FR/EN généré depuis 504 concepts. &
Conserve le statut de provenance et les codes protégés.\\
\hline
\code{src/language-model.js} &
Prédiction, score de fluidité et classement hybrides. &
Ne peut pas autoriser seul une traduction.\\
\hline
\texttt{src/neural-\allowbreak language-model.js} &
Inférence JavaScript des poids int8. &
Pas de bibliothèque d’inférence ni de chargement réseau.\\
\hline
\code{chrome.storage.local} &
Mémorise activation et langue choisie. &
Aucun texte de page, montant, nom ou historique.\\
\hline
\end{longtable}

\subsection{Séquence d’une activation}

\begin{figure}[H]
\centering
\begin{tikzpicture}[
  actor/.style={draw=Teal,fill=Mint,rounded corners=1.5mm,
    text width=22mm,minimum height=10mm,align=center,font=\small\bfseries},
  msg/.style={-{Latex[length=2mm]},line width=0.7pt,color=TealDark},
  life/.style={dashed,color=Line,thick}
]
\node[actor] (u) {Utilisateur};
\node[actor,right=11mm of u] (p) {Popup};
\node[actor,right=11mm of p] (c) {Content script};
\node[actor,right=11mm of c] (e) {Moteur};
\node[actor,right=11mm of e] (d) {Site};
\foreach \n in {u,p,c,e,d}
  \draw[life] (\n.south) -- ++(0,-78mm);
\draw[msg] ($(u.south)+(0,-8mm)$) -- node[above,font=\scriptsize]{choisit Fulfulde} ($(p.south)+(0,-8mm)$);
\draw[msg] ($(p.south)+(0,-18mm)$) -- node[above,font=\scriptsize]{active KORA} ($(c.south)+(0,-18mm)$);
\draw[msg] ($(c.south)+(0,-28mm)$) -- node[above,font=\scriptsize]{lit DOM et valeurs} ($(d.south)+(0,-28mm)$);
\draw[msg] ($(c.south)+(0,-38mm)$) -- node[above,font=\scriptsize]{segment protégé} ($(e.south)+(0,-38mm)$);
\draw[msg] ($(e.south)+(0,-48mm)$) -- node[above,font=\scriptsize]{sens et intégrité} ($(c.south)+(0,-48mm)$);
\draw[msg] ($(c.south)+(0,-58mm)$) -- node[above,font=\scriptsize]{texte visible traduit} ($(d.south)+(0,-58mm)$);
\draw[msg] ($(d.south)+(0,-68mm)$) -- node[above,font=\scriptsize]{nouveau contenu} ($(c.south)+(0,-68mm)$);
\end{tikzpicture}
\caption{La popup transmet le choix de langue; le content script lit et rend
le texte; le moteur contrôle le sens et l’intégrité sans modifier la logique du site.}
\end{figure}

\subsection{Pipeline de traduction contrôlée}

\begin{figure}[H]
\centering
\resizebox{0.97\linewidth}{!}{%
\begin{tikzpicture}[
  node distance=5mm and 7mm,
  action/.style={draw=Line,fill=Paper,rounded corners=1.5mm,
    text width=28mm,minimum height=11mm,align=center,font=\small},
  decision/.style={diamond,draw=Teal,fill=Mint,aspect=2.3,
    text width=25mm,align=center,font=\small,inner sep=1.5pt},
  safe/.style={action,draw=Teal,fill=Mint},
  stop/.style={action,draw=Gold,fill=Warm},
  arrow/.style={-{Latex[length=2mm]},thick,color=TealDark}
]
\node[action] (source) {Segment FR ou EN};
\node[decision,right=of source] (protected) {Valeur ou code protégé ?};
\node[action,right=of protected] (mask) {Masquer montants, dates, comptes et références};
\node[decision,right=of mask] (known) {Mémoire, faute unique, gabarit ou glossaire ?};
\node[safe,above right=8mm and 9mm of known] (candidate) {Candidat \code{fub} dont le sens est couvert};
\node[stop,below right=8mm and 9mm of known] (original) {Conserver l’original};
\node[decision,right=10mm of candidate] (coverage) {Couverture suffisante ?};
\node[action,right=of coverage] (rank) {Classement de fluidité};
\node[action,right=of rank] (restore) {Restaurer les valeurs};
\node[decision,right=of restore] (integrity) {Occurrences identiques ?};
\node[safe,right=of integrity] (render) {Rendre le texte};

\draw[arrow] (source) -- (protected);
\draw[arrow] (protected) -- node[above,font=\scriptsize]{non} (mask);
\draw[arrow] (protected) |- node[pos=.25,left,font=\scriptsize]{oui} (original);
\draw[arrow] (mask) -- (known);
\draw[arrow] (known) -- node[above left,font=\scriptsize]{oui} (candidate);
\draw[arrow] (known) -- node[below left,font=\scriptsize]{non} (original);
\draw[arrow] (candidate) -- (coverage);
\draw[arrow] (coverage) -- node[above,font=\scriptsize]{oui} (rank);
\draw[arrow] (coverage) |- node[pos=.25,right,font=\scriptsize]{non} (original);
\draw[arrow] (rank) -- (restore);
\draw[arrow] (restore) -- (integrity);
\draw[arrow] (integrity) -- node[above,font=\scriptsize]{oui} (render);
\draw[arrow] (integrity) |- node[pos=.25,right,font=\scriptsize]{non} (original);
\end{tikzpicture}}
\caption{La fluidité intervient après la détermination du sens et avant le
dernier contrôle d’intégrité.}
\end{figure}

Ordre de décision : mémoire Safinum exacte; mémoire web/finance générale;
terminologie contrôlée; correction orthographique non ambiguë; gabarit typé;
phrasebook par correspondance longue; abstention. Pour un message financier
critique, la couverture requise est de 100\,\%; elle est de 80\,\% pour un
segment non critique, sous réserve du verrou d’intégrité.

\subsection{Contenu dynamique}

\begin{figure}[H]
\centering
\resizebox{\linewidth}{!}{%
\begin{tikzpicture}[
  node distance=4.5mm,
  flow/.style={draw=Line,fill=Paper,rounded corners=1.5mm,
    minimum height=10mm,text width=29mm,align=center,font=\small},
  arrow/.style={-{Latex[length=2mm]},thick,color=TealDark}
]
\node[flow] (mut) {Mutation DOM ou événement utilisateur};
\node[flow,right=of mut] (deb) {Regroupement court des changements};
\node[flow,right=of deb] (disc) {Découverte du document et des racines ouvertes};
\node[flow,right=of disc] (class) {Classement : source ou valeur protégée};
\node[flow,right=of class] (trans) {Même pipeline de traduction};
\node[flow,right=of trans] (apply) {Mise à jour idempotente du texte};
\draw[arrow] (mut)--(deb);
\draw[arrow] (deb)--(disc);
\draw[arrow] (disc)--(class);
\draw[arrow] (class)--(trans);
\draw[arrow] (trans)--(apply);
\end{tikzpicture}}
\caption{Traitement des contenus ajoutés après le chargement de la page.}
\end{figure}

Le content script combine \code{MutationObserver}, les événements
\code{click}/\code{change}/\code{input} et un rescan de sûreté toutes les
500\,ms lorsque l’onglet est visible. Il prend en charge les écrans montés par
framework, les frais recalculés, les erreurs, les attributs et les options,
sans remplacer le composant d’origine.

\subsection{Construction et artefacts}

\begin{figure}[H]
\centering
\resizebox{\linewidth}{!}{%
\begin{tikzpicture}[
  node distance=5mm and 7mm,
  b/.style={draw=Line,fill=Paper,rounded corners=1.5mm,
    minimum height=10mm,text width=27mm,align=center,font=\small},
  artifact/.style={b,draw=Teal,fill=Mint},
  arrow/.style={-{Latex[length=2mm]},thick,color=TealDark}
]
\node[b] (src) {Sources JS, données et manifest};
\node[b,right=of src] (term) {Compilation terminologique};
\node[b,right=of term] (bundle) {Bundle esbuild + ressources};
\node[artifact,right=of bundle] (unpacked) {\code{kora-unpacked/}};
\node[artifact,above right=5mm and 7mm of unpacked] (zip) {{\ttfamily\scriptsize kora-\\extension.zip}};
\node[artifact,below right=5mm and 7mm of unpacked] (json) {Modèle JSON exportable};
\draw[arrow] (src)--(term);
\draw[arrow] (term)--(bundle);
\draw[arrow] (bundle)--(unpacked);
\draw[arrow] (unpacked)--(zip);
\draw[arrow] (bundle)|-(json);
\end{tikzpicture}}
\caption{Chaîne de build : le packaging échoue si l’archive dépasse 10 Mio.}
\end{figure}

Le ZIP final ne contient ni corpus brut, ni PyTorch, ni Puppeteer, ni clé
secrète. Il contient le manifest, la popup, les icônes et les bundles où sont
déjà intégrés moteur, mémoires et poids KORA.
